\documentclass[12pt,a4paper]{ctexart}

% ============================================================
%                       宏包设置
% ============================================================

\usepackage{amsmath, amssymb, amsfonts}
\usepackage{geometry}
\usepackage{fancyhdr}
\usepackage{graphicx}
\usepackage{xcolor}
\usepackage{setspace}
\usepackage[most]{tcolorbox}
\usepackage{enumitem}

% ============================================================
%                 清新蓝绿色配色
% ============================================================

% 主色：清爽青绿色
\definecolor{SZURed}{HTML}{2F9E8F}

% 题目区域：浅薄荷绿
\definecolor{SZULightRed}{HTML}{EAF7F4}

% 答案区域：浅天空蓝
\definecolor{SZULightRose}{HTML}{EEF7FC}

% 边框：柔和蓝绿色
\definecolor{SZUSoftBorder}{HTML}{9CCFC8}

% 强调文字：沉稳蓝绿色
\definecolor{SZUTextRed}{HTML}{256B72}

% ============================================================
%                       页面设置
% ============================================================

\geometry{
    left=2.5cm,
    right=2.5cm,
    top=2.5cm,
    bottom=2.5cm
}

\setlength{\parindent}{0pt}
\setlength{\parskip}{0.6em}
\linespread{1.15}

% ============================================================
%                       课程信息
% ============================================================

\newcommand{\coursename}{概率论与数理统计}
\newcommand{\department}{深圳大学生物医学工程}
\newcommand{\homeworknumber}{作业 2}

% ============================================================
%                  学生信息 —— 请学生修改
% ============================================================

\newcommand{\studentname}{ }
\newcommand{\studentid}{ }

% ============================================================
%                       页眉页脚
% ============================================================

\pagestyle{fancy}
\fancyhf{}

\lhead{\textcolor{SZUTextRed}{\coursename}}
\rhead{\textcolor{SZUTextRed}{\homeworknumber}}
\cfoot{\textcolor{SZUTextRed}{\thepage}}

\renewcommand{\headrulewidth}{0.4pt}
\renewcommand{\headrule}{
    \hbox to\headwidth{
        \color{SZUSoftBorder}\leaders\hrule height \headrulewidth\hfill
    }
}

\setlength{\headheight}{15pt}

% ============================================================
%                    题目编号与环境设置
% ============================================================

\newcounter{problem}

% -------------------------
% 题目环境
% -------------------------

\newenvironment{problem}[1][]
{
    \stepcounter{problem}
    \begin{tcolorbox}[
        title={题目 \theproblem\if&#1&\relax\else\quad[#1]\fi},
        colback=SZULightRed,
        colframe=SZURed,
        coltitle=white,
        fonttitle=\bfseries,
        boxrule=0.8pt,
        arc=2mm,
        left=3mm,
        right=3mm,
        top=2.5mm,
        bottom=2.5mm,
        breakable
    ]
}
{
    \end{tcolorbox}
}

% -------------------------
% 答案环境
% -------------------------

\newenvironment{answer}
{
    \begin{tcolorbox}[
        title={学生答案},
        colback=SZULightRose,
        colframe=SZUSoftBorder,
        coltitle=SZUTextRed,
        fonttitle=\bfseries,
        boxrule=0.7pt,
        arc=2mm,
        left=3mm,
        right=3mm,
        top=2.5mm,
        bottom=2.5mm,
        breakable
    ]
}
{
    \end{tcolorbox}
}

% ============================================================
%                         正文
% ============================================================

\begin{document}

% ============================================================
%                         标题
% ============================================================

\begin{center}

    {\LARGE\bfseries\textcolor{SZURed}{\coursename}}

    \vspace{0.4em}

    {\large \department}

    \vspace{1em}

    {\Large\bfseries\textcolor{SZUTextRed}{\homeworknumber}}

\end{center}

\vspace{1em}

% ============================================================
%                       学生信息
% ============================================================

\begin{center}

\begin{tabular}{rl@{\hspace{3cm}}rl}

\textcolor{SZUTextRed}{\bfseries 姓名：}
&
\studentname
&
\textcolor{SZUTextRed}{\bfseries 学号：}
&
\studentid

\end{tabular}

\end{center}

\vspace{0.7em}

{\color{SZUSoftBorder}\hrule}

\vspace{1.5em}



% ============================================================
% ============================================================
%
%                           题目 1
%
% ============================================================
% ============================================================

\begin{problem}[课本习题3]

设在15只同类型的零件中有2只是次品，在其中取3次，每次任取1只，作不放回抽样。以$X$表示取出的次品的只数。
\begin{enumerate}[label=(\arabic*)]
    \item 求$X$的分布律。
    \item 画出分布律的图形。
\end{enumerate}

\end{problem}


% ============================================================
%
%              ↓↓↓↓↓↓ 学生答案区域 ↓↓↓↓↓↓
%
%          请只在 answer 环境中填写你的答案
%
% ============================================================

\begin{answer}

% ==========================================================
%                 请从下一行开始填写答案
% ==========================================================






% ==========================================================
%                 请在上一行结束填写答案
% ==========================================================

\end{answer}

% ============================================================
%
%              ↑↑↑↑↑↑ 学生答案区域 ↑↑↑↑↑↑
%
% ============================================================


\vspace{1em}



% ============================================================
% ============================================================
%
%                           题目 2
%
% ============================================================
% ============================================================

\begin{problem}[课本习题9]

有一大批产品，其验收方案如下，先作第一次检验：从中任取10件，经检验无次品接受这批产品，次品数大于2拒收；否则作第二次检验，其做法是从中再任取5件，仅当5件中无次品时接受这批产品。若产品的次品率为10\%，求
\begin{enumerate}[label=(\arabic*)]
    \item 这批产品经第一次检验就能接受的概率。
    \item 需作第二次检验的概率。
    \item 这批产品按第二次检验的标准被接受的概率。
    \item 这批产品在第一次检验未能作决定且第二次检验时被通过的概率。
    \item 这批产品被接受的概率。
\end{enumerate}
\end{problem}


% ============================================================
%
%              ↓↓↓↓↓↓ 学生答案区域 ↓↓↓↓↓↓
%
% ============================================================

\begin{answer}

% ==========================================================
%                 请从下一行开始填写答案
% ==========================================================





% ==========================================================
%                 请在上一行结束填写答案
% ==========================================================

\end{answer}

% ============================================================
%
%              ↑↑↑↑↑↑ 学生答案区域 ↑↑↑↑↑↑
%
% ============================================================


\vspace{1em}



% ============================================================
% ============================================================
%
%                           题目 3
%
% ============================================================
% ============================================================

\begin{problem}[课本习题12]

一电话总机每分钟收到呼唤的次数服从参数为4的泊松分布。求
\begin{enumerate}[label=(\arabic*)]
    \item 某一分钟恰有8次呼唤的概率；
    \item 某一分钟的呼唤次数大于3的概率。
\end{enumerate}

\end{problem}


% ============================================================
%
%              ↓↓↓↓↓↓ 学生答案区域 ↓↓↓↓↓↓
%
% ============================================================

\begin{answer}

% ==========================================================
%                 请从下一行开始填写答案
% ==========================================================





% ==========================================================
%                 请在上一行结束填写答案
% ==========================================================

\end{answer}

% ============================================================
%
%              ↑↑↑↑↑↑ 学生答案区域 ↑↑↑↑↑↑
%
% ============================================================


\vspace{1em}



% ============================================================
% ============================================================
%
%                           题目 4
%
% ============================================================
% ============================================================

\begin{problem}[课本习题14]

某人家中在时间间隔$t$（以h计）内接到电话的次数$X$服从参数为$2t$的泊松分布。
\begin{enumerate}[label=(\arabic*)]
    \item 若他外出计划用时10 min，问其间有电话铃响一次的概率是多少？
    \item 若他希望外出时没有电话的概率至少为0.5，问他外出应控制最长时间是多少？
\end{enumerate}

\end{problem}


% ============================================================
%
%              ↓↓↓↓↓↓ 学生答案区域 ↓↓↓↓↓↓
%
% ============================================================

\begin{answer}

% ==========================================================
%                 请从下一行开始填写答案
% ==========================================================





% ==========================================================
%                 请在上一行结束填写答案
% ==========================================================

\end{answer}

% ============================================================
%
%              ↑↑↑↑↑↑ 学生答案区域 ↑↑↑↑↑↑
%
% ============================================================


\vspace{1em}



% ============================================================
% ============================================================
%
%                           题目 5
%
% ============================================================
% ============================================================

\begin{problem}[课本习题19]

以$X$表示某商店从早晨开始营业起直到第一个顾客到达的等待时间（以min计），$X$的分布函数是
\[
F_X(x)=
\begin{cases}
1-\mathrm{e}^{-0.4x}, & x>0,\\
0, & x\le 0.
\end{cases}
\]
求下述概率：
\begin{enumerate}[label=(\arabic*)]
    \item $P\{\text{至多 }3\ \text{min}\}$。
    \item $P\{\text{至少 }4\ \text{min}\}$。
    \item $P\{3\ \text{min 至 }4\ \text{min 之间}\}$。
    \item $P\{\text{至多 }3\ \text{min 或至少 }4\ \text{min}\}$。
    \item $P\{\text{恰好 }2.5\ \text{min}\}$。
\end{enumerate}

\end{problem}


% ============================================================
%
%              ↓↓↓↓↓↓ 学生答案区域 ↓↓↓↓↓↓
%
% ============================================================

\begin{answer}

% ==========================================================
%                 请从下一行开始填写答案
% ==========================================================





% ==========================================================
%                 请在上一行结束填写答案
% ==========================================================

\end{answer}

% ============================================================
%
%              ↑↑↑↑↑↑ 学生答案区域 ↑↑↑↑↑↑
%
% ============================================================


\vspace{1em}



% ============================================================
% ============================================================
%
%                           题目 6
%
% ============================================================
% ============================================================

\begin{problem}[课本习题23]

某种型号器件的寿命$X$（以小时计）具有概率密度
\[
f(x)=
\begin{cases}
\dfrac{1000}{x^2}, & x>1000,\\
0, & \text{其他}.
\end{cases}
\]
现有一大批此种器件（设器件损坏与否相互独立），任取5只，问其中至少有2只寿命大于1500小时的概率是多少？

\end{problem}


% ============================================================
%
%              ↓↓↓↓↓↓ 学生答案区域 ↓↓↓↓↓↓
%
% ============================================================

\begin{answer}

% ==========================================================
%                 请从下一行开始填写答案
% ==========================================================





% ==========================================================
%                 请在上一行结束填写答案
% ==========================================================

\end{answer}

% ============================================================
%
%              ↑↑↑↑↑↑ 学生答案区域 ↑↑↑↑↑↑
%
% ============================================================


\vspace{1em}



% ============================================================
% ============================================================
%
%                           题目 7
%
% ============================================================
% ============================================================

\begin{problem}[课本习题26]

设$X\sim N(3,2^2)$。
\begin{enumerate}[label=(\arabic*)]
    \item 求$P\{2<X\le 5\},P\{-4<X\le 10\},P\{|X|>2\},P\{X>3\}$。
    \item 确定$c$，使得$P\{X>c\}=P\{X\le c\}$。
    \item 设$d$满足$P\{X>d\}\ge 0.9$，问$d$至多为多少？
\end{enumerate}

\end{problem}


% ============================================================
%
%              ↓↓↓↓↓↓ 学生答案区域 ↓↓↓↓↓↓
%
% ============================================================

\begin{answer}

% ==========================================================
%                 请从下一行开始填写答案
% ==========================================================





% ==========================================================
%                 请在上一行结束填写答案
% ==========================================================

\end{answer}

% ============================================================
%
%              ↑↑↑↑↑↑ 学生答案区域 ↑↑↑↑↑↑
%
% ============================================================


\vspace{1em}



% ============================================================
% ============================================================
%
%                           题目 8
%
% ============================================================
% ============================================================

\begin{problem}[课本习题30]

设在一电路中，电阻两端的电压(V)服从$N(120,2^2)$，今独立测量了5次，试确定有2次测定值落在区间$[118,122]$之外的概率。
\end{problem}


% ============================================================
%
%              ↓↓↓↓↓↓ 学生答案区域 ↓↓↓↓↓↓
%
% ============================================================

\begin{answer}

% ==========================================================
%                 请从下一行开始填写答案
% ==========================================================





% ==========================================================
%                 请在上一行结束填写答案
% ==========================================================

\end{answer}

% ============================================================
%
%              ↑↑↑↑↑↑ 学生答案区域 ↑↑↑↑↑↑
%
% ============================================================


\vspace{1em}



% ============================================================
% ============================================================
%
%                           题目 9
%
% ============================================================
% ============================================================

\begin{problem}[课本习题34]

设随机变量$X$在区间$(0,1)$服从均匀分布。
\begin{enumerate}[label=(\arabic*)]
    \item 求$Y=\mathrm{e}^X$的概率密度。
    \item 求$Y=-2\ln X$的概率密度。
\end{enumerate}

\end{problem}


% ============================================================
%
%              ↓↓↓↓↓↓ 学生答案区域 ↓↓↓↓↓↓
%
% ============================================================

\begin{answer}

% ==========================================================
%                 请从下一行开始填写答案
% ==========================================================





% ==========================================================
%                 请在上一行结束填写答案
% ==========================================================

\end{answer}

% ============================================================
%
%              ↑↑↑↑↑↑ 学生答案区域 ↑↑↑↑↑↑
%
% ============================================================


\vspace{1em}



% ============================================================
% ============================================================
%
%                           题目 10
%
% ============================================================
% ============================================================

\begin{problem}[课本习题37]

设随机变量$X$的概率密度为
\[
f(x)=
\begin{cases}
\dfrac{2x}{\pi^2}, & 0<x<\pi,\\
0, & \text{其他}.
\end{cases}
\]
求$Y=\sin X$的概率密度。

\end{problem}


% ============================================================
%
%              ↓↓↓↓↓↓ 学生答案区域 ↓↓↓↓↓↓
%
% ============================================================

\begin{answer}

% ==========================================================
%                 请从下一行开始填写答案
% ==========================================================





% ==========================================================
%                 请在上一行结束填写答案
% ==========================================================

\end{answer}

% ============================================================
%
%              ↑↑↑↑↑↑ 学生答案区域 ↑↑↑↑↑↑
%
% ============================================================




\end{document}
```
